\documentclass[12pt]{article}
\usepackage[utf8]{inputenc}
\usepackage[T2A]{fontenc}
\usepackage[russian]{babel}
\usepackage{amsmath, amssymb}
\usepackage{geometry}
\usepackage{graphicx}
\usepackage{titlesec} % Для настройки заголовков
\usepackage{enumitem} % Для настройки списков
\usepackage{fancyhdr} % Для колонтитулов
\usepackage{tabularx} % Для таблиц
\usepackage{lastpage} % Для номера последней страницы
\usepackage{hyperref}


% --- НАСТРОЙКА СТИЛЯ ---
\geometry{a4paper, left=25mm, right=15mm, top=20mm, bottom=20mm, headheight=15pt}
\setlength{\parindent}{0pt} % Убираем абзацный отступ
\setlist[enumerate]{topsep=5pt, itemsep=8pt, leftmargin=25pt} % Настройка списков
\titleformat{\section}{\large\bfseries\centering}{\thesection}{1em}{} % Стиль секций

% --- КОЛОНТИТУЛЫ (со 2-й страницы) ---
\pagestyle{fancy}
\fancyhf{}
\fancyhead[L]{
\includegraphics[height=1.1cm]{logo.png} \textbf{MathCore} | MOCK Test-2}
\fancyhead[R]{Abdunazarov M.X.}
\fancyfoot[C]{}
\fancyfoot[R]{%

\includegraphics[height=1cm]{tg.png}%
\href{https://t.me/Abd_mardon}{\textbf{@Abd\_mardon}}%
}
\renewcommand{\headrulewidth}{0.4pt}

\title{MOCK Test-2}
\author{Abdunazarov M.X.}
\date{} 

\begin{document}

\maketitle
\begin{center}
    
\includegraphics[width=0.5\linewidth]{logo.png}
\end{center}
\begin{enumerate}

\item В равностороннем треугольнике со стороной 4 случайным образом выбирается точка. Внутри треугольника построена вписанная окружность. Найдите вероятность того, что выбранная точка попадёт внутрь окружности.
 
\textbf{}
\begin{enumerate}
    \item[A)] \(\frac{\pi\sqrt{3}}{6}\)
    \item[B)] \(\frac{\pi\sqrt{3}}{9}\)
    \item[C)] 0.4
    \item[D)] \(\frac{2}{\pi}\)
\end{enumerate}

\vspace{0.cm} % Пустое место для решения

\item Вычислите: \(\frac{15}{14} + \frac{1515}{1414} + \frac{151515}{141414} + \cdots + \frac{15151515151515}{14141414141414}\) \\
\textbf{}
\begin{enumerate}
    \item[A)] 15
    \item[B)] 3
    \item[C)] 7,5
    \item[D)] 8,5
\end{enumerate}

\vspace{0.cm} % Пустое место для решения

\item \( a, b, c \) — положительные целые числа, удовлетворяющие условиям \( \text{НОД}(b; c) = 5 \); \( \text{НОД}(c; d) = 7 \). Найдите наименьшее значение \( b + c + d \). \\
\textbf{}
\begin{enumerate}
    \item[A)] 12
    \item[B)] 24
    \item[C)] 38
    \item[D)] 47
\end{enumerate}

\vspace{0.cm} % Пустое место для решения

\item Найдите произведение действительных корней уравнения \((x - 1)^4 + 2x = x^2 + 73\). \\
\textbf{}
\begin{enumerate}
    \item[A)] 5
    \item[B)] \(-7\)
    \item[C)] 4
    \item[D)] \(-8\)
\end{enumerate}

\vspace{0.cm} % Пустое место для решения

\item Упростите выражение:
\[
\frac{2x+y-\frac{4y^2-x^2}{y}}{x^3+2x^2y-3xy^2} + \left( \frac{y^3x-2x^2y^2+x^3y}{y^2-x^2} \right)^{-1}
\]
\textbf{}
\begin{enumerate}
    \item[A)] \(\frac{-2}{x(y-x)}\)
    \item[B)] \(\frac{-2y}{x(y-x)}\)
    \item[C)] \(\frac{2y}{x(y-x)}\)
    \item[D)] \(\frac{2}{x(y-x)}\)
\end{enumerate}

\vspace{0.cm} % Пустое место для решения

\item Жасур задумал одно положительное целое число. К его записи справа приписали цифру 5, и из полученного числа вычли квадрат задуманного числа. Разность разделили на задуманное число, затем из результата вычли задуманное число. В итоге получилась 1. На какое из следующих чисел делится задуманное число без остатка? \\
\textbf{}
\begin{enumerate}
    \item[A)] 5
    \item[B)] 7
    \item[C)] 9
    \item[D)] 4
\end{enumerate}

\vspace{0.cm} % Пустое место для решения

\item Сосуд, заполненный 12\%-м раствором соли, опустошили на 1 литр раствора и долили водой, затем снова отлили 1 литр раствора и долили водой. В результате в сосуде получился 3\%-й раствор соли. Найдите вместимость сосуда. \\
\textbf{}
\begin{enumerate}
    \item[A)] 2 л
    \item[B)] 5 л
    \item[C)] 4 л
    \item[D)] 1,5 л
\end{enumerate}

\vspace{0.cm} % Пустое место для решения

\item В автогонках из двух этапов, три автомобиля двигались с постоянной скоростью. Известно следующее: Автомобиль I на первом этапе двигался со скоростью 150 км/ч и завершил первый этап за 12 минут. Автомобиль II прошел всю гонку со скоростью 120 км/ч и завершил ее за 20 минут. Автомобиль III завершил второй этап за 6 минут. Найдите его скорость (в км/ч). \\
\textbf{}
\begin{enumerate}
    \item[A)] 80
    \item[B)] 90
    \item[C)] 100
    \item[D)] 110
\end{enumerate}

\vspace{0.cm} % Пустое место для решения

\item Найдите сумму наименьшего положительного целого и наибольшего отрицательного целого значений параметра \(a\), при которых уравнение \(x^2 + 7|x| + 12 - a - a^2 = 0\) имеет два различных действительных корня. \\
\textbf{}
\begin{enumerate}
    \item[A)] 4
    \item[B)] 5
    \item[C)] \(-1\)
    \item[D)] 0
\end{enumerate}

\vspace{0.cm} % Пустое место для решения

\item На основе графиков показательных функций \(y = a^x\), \(y = b^x\), \(y = c^x\) упростите выражение \( |b - a| + |c - a| - |c - b| \). \\
\begin{center}
    
        \centering
        \includegraphics[width=0.5\linewidth]{1.png}
    
\end{center}
\textbf{}
\begin{enumerate}
    \item[A)] \(2a - 2b\)
    \item[B)] \(2b - 2a\)
    \item[C)] \(a - c\)
    \item[D)] 0
\end{enumerate}

\vspace{0.cm} % Пустое место для решения

\item При каких значениях параметра \(m\) показательная функция \(y = (m - 2)^x\) является возрастающей, а логарифмическая функция \(y = \log_{2m-7}x\) — убывающей? \\
\textbf{}
\begin{enumerate}
    \item[A)] \((3; 4)\)
    \item[B)] \((3; +\infty)\)
    \item[C)] \((3.5; 4)\)
    \item[D)] \((4; +\infty)\)
\end{enumerate}

\vspace{0.cm} % Пустое место для решения

\item Найдите сумму корней уравнения \(\sin 3x - \sin 7x = \sqrt{3} \sin 2x\) на промежутке \([0; \frac{\pi}{2}]\). \\
\textbf{}
\begin{enumerate}
    \item[A)] \(\frac{9\pi}{10}\)
    \item[B)] \(\frac{2\pi}{3}\)
    \item[C)] \(\frac{5\pi}{6}\)
    \item[D)] \(\frac{2\pi}{5}\)
\end{enumerate}

\vspace{0.cm} % Пустое место для решения

\item Какое наименьшее число нужно прибавить к корню уравнения, чтобы получить квадрат некоторого числа?
\[
\frac{6}{-\sqrt{x} + \sqrt{x + 6}} + \frac{3}{\sqrt{x - 3} - \sqrt{x}} = 1
\]
\textbf{}
\begin{enumerate}
    \item[A)] 4
    \item[B)] 5
    \item[C)] 6
    \item[D)] 7
\end{enumerate}

\vspace{0.cm} % Пустое место для решения

\item Для ненулевых действительных чисел \(a\) и \(b\), числа \(a^2b\), \(3ab\), \(a + 2b\) в указанном порядке образуют одновременно арифметическую и геометрическую прогрессии. Найдите значение \(b\). \\
\textbf{}
\begin{enumerate}
    \item[A)] \(\frac{3}{7}\)
    \item[B)] \(\frac{3}{5}\)
    \item[C)] \(\frac{3}{2}\)
    \item[D)] \(\frac{7}{3}\)
\end{enumerate}

\vspace{0.cm} % Пустое место для решения

\item В информационно-ресурсном центре 30 компьютеров соединены друг с другом кабелями. Если сумма входящих и исходящих кабелей для каждого компьютера равна 6, сколько всего кабелей потребуется? \\
\textbf{}
\begin{enumerate}
    \item[A)] 90
    \item[B)] 40
    \item[C)] 180
    \item[D)] 200
\end{enumerate}

\vspace{0.cm} % Пустое место для решения

\item В первой урне 3 красных и 2 синих шара, во второй — 4 синих и 2 красных шара. Из первой урны берут шар и перекладывают во вторую. Затем из второй урны берут шар. Найдите вероятность того, что шар, взятый из второй урны, будет того же цвета, что и шар, взятый из первой урны. \\
\textbf{}
\begin{enumerate}
    \item[A)] \(\frac{2}{7}\)
    \item[B)] \(\frac{19}{35}\)
    \item[C)] \(\frac{9}{35}\)
    \item[D)] \(\frac{3}{5}\)
\end{enumerate}

\vspace{0.cm} % Пустое место для решения

\item Если \(x^2 + x + 1 = 0\), найдите числовое значение выражения \(\left(x^3 + \frac{1}{x^3}\right)^3\). \\
\textbf{}
\begin{enumerate}
    \item[A)] 8
    \item[B)] 1
    \item[C)] \(-1\)
    \item[D)] \(-8\)
\end{enumerate}

\vspace{0.cm} % Пустое место для решения

\item На основе данных на рисунке найдите наибольшую площадь прямоугольника.
\begin{center}
   
       \centering
       \includegraphics[width=0.5\linewidth]{2.png}
   
\end{center}
\textbf{}
\begin{enumerate}
    \item[A)] 48
    \item[B)] 36
    \item[C)] 28
    \item[D)] 54
\end{enumerate}

\vspace{0.cm} % Пустое место для решения

\item На рисунке изображен график производной функции \(f(x)\) на интервале (\(-6; 12\)). Найдите отношение суммы длин промежутков возрастания функции к сумме длин промежутков убывания на интервале (\(-6; 12\)).
\begin{center}
   
       \centering
       \includegraphics[width=0.5\linewidth]{3.png}
   
\end{center}
\textbf{}
\begin{enumerate}
    \item[A)] 2:1
    \item[B)] \(\frac{1}{2}\)
    \item[C)] \(\frac{11}{7}\)
    \item[D)] \(\frac{13}{5}\)
\end{enumerate}

\vspace{0.cm} % Пустое место для решения

\item Вычислите интеграл: 
\[
\int \frac{x}{\cos^2 (x^2 + 1)}  dx
\]
\textbf{}
\begin{enumerate}
    \item[A)] \(\frac{1}{2} \tg(x^2 + 1) + C\)
    \item[B)] \(2 \tg(x^2 + 1) + C\)
    \item[C)] \(\frac{1}{2} \ctg(x^2 + 1) + C\)
    \item[D)] \(\ctg(x^2 + 1) + C\)
\end{enumerate}



\item \(S_1, S_2, S_3\) обозначают площади соответствующих областей. Основываясь на этом, найдите значение отношения \(\frac{S_1 + S_2}{S_3}\).
\begin{center}
 
       \centering
       \includegraphics[width=0.4\linewidth]{4.png}
  
\end{center}
\textbf{}
\begin{enumerate}
    \item[A)] 3:1
    \item[B)] 2:1
    \item[C)] 3:2
    \item[D)] 1
\end{enumerate}

\vspace{0.cm} % Пустое место для решения

\item Сопоставьте следующие утверждения?

     
        1. Периметр прямоугольного треугольника равен \(P = 2(c + r)\), где \(c\) — гипотенуза, \(r\) — радиус вписанной окружности.\\
       
        2. Центр окружности, описанной около правильного треугольника, лежит на его биссектрисе. \\
        
        3. Четырехугольник, полученный последовательным соединением середин сторон произвольного четырехугольника, является параллелограммом. \\
        
 
\textbf{}
\begin{enumerate}
    \item[A)] 1-Верно, 2-Верно, 3-Неверно
    \item[B)] 1-Неверно, 2-Верно, 3-Верно
    \item[C)] 1-Неверно, 2-Неверно, 3-Неверно
    \item[D)] 1-Верно, 2-Верно, 3-Верно
\end{enumerate}

\vspace{0.cm} % Пустое место для решения


\item Окружность разделена на три дуги в отношении 13 : 14 : 9, и точки деления соединены, образуя треугольник. Если меньшая сторона полученного треугольника равна \(\sqrt[4]{12}\) см, найдите радиус этой окружности. \\
\textbf{}
\begin{enumerate}
    \item[A)] \(\sqrt[4]{4}\)
    \item[B)] 4,5
    \item[C)] 9
    \item[D)] \(\sqrt[4]{3}\)
\end{enumerate}

\vspace{0.cm} % Пустое место для решения

\item В параллелограмме \(ABCD\) проведена высота \(DE\). Точки \(E\) и \(F\) являются серединами сторон \(AB\) и \(DE\) соответственно. Если площадь параллелограмма \(ABCD\) равна 48, найдите площадь треугольника \(CGF\). \\
\begin{center}
    \centering
    \includegraphics[width=0.5\linewidth]{5.png}
\end{center}
\textbf{}
\begin{enumerate}
    \item[A)] 4
    \item[B)] 2
    \item[C)] 8
    \item[D)] 6
\end{enumerate}

\vspace{0.cm} % Пустое место для решения

\item В трапеции \(ABCD\) \(AE = BE\). Если \(S_{CKDE} = 9\) и \(S_{CDK} = 3\), найдите площадь трапеции. \\
\begin{center}
    \centering
    \includegraphics[width=0.5\linewidth]{6.png}
\end{center}
\textbf{}
\begin{enumerate}
    \item[A)] 54
    \item[B)] 52
    \item[C)] 48
    \item[D)] 44
\end{enumerate}

\vspace{0.cm} % Пустое место для решения

\item \(ABCDEFGH\) — правильный восьмиугольник. Если \(\overrightarrow{AH} = \vec{x}\), \(\overrightarrow{AE} = \vec{y}\), выразите вектор \(\overrightarrow{DC} + \overrightarrow{CB} + \overrightarrow{BA}\) через векторы \(\vec{x}\) и \(\vec{y}\). \\
\textbf{}
\begin{enumerate}
    \item[A)] \(\vec{x} + \vec{y}\)
    \item[B)] \(\vec{x} - \vec{y}\)
    \item[C)] \(\vec{y} - \vec{x}\)
    \item[D)] \(-\vec{x} - \vec{y}\)
\end{enumerate}

\vspace{0.cm} % Пустое место для решения

\item Дан прямоугольный треугольник \(ABC\). Точки \(K\) и \(M\) являются центрами масс треугольников \(ABD\) и \(BDH\) соответственно. Найдите длину отрезка \(|KM|\). \\
\begin{center}
    \centering
    \includegraphics[width=0.5\linewidth]{7.png}
\end{center}
\textbf{}
\begin{enumerate}
    \item[A)] \(2\sqrt{3}\)
    \item[B)] \(2\sqrt{2}\)
    \item[C)] 3
    \item[D)] 4
\end{enumerate}

\vspace{0.cm} % Пустое место для решения

\item Даны четырехугольники \(ABCD\) и \(AEFK\). Если \(BC = FL = 4\) см, \(LC = 3\) см, \(AB = 10\) см и \(KF = 8\) см, найдите \(AK = x\) из приведенного чертежа. \\
\begin{center}
    \centering
    \includegraphics[width=0.5\linewidth]{11.png}
\end{center}
\textbf{}
\begin{enumerate}
    \item[A)] \(3\sqrt{3}\)
    \item[B)] \(4\sqrt{2}\)
    \item[C)] 5
    \item[D)] 6
\end{enumerate}

\vspace{0.cm} % Пустое место для решения

\item В треугольник с основанием \(a\) и высотой \(h\), опущенной на это основание, вписан параллелограмм так, что одна из его сторон лежит на основании \(a\). Как нужно выбрать основание параллелограмма, чтобы его площадь была максимальной? \\
\textbf{}
\begin{enumerate}
    \item[A)] \(\frac{a}{2}\)
    \item[B)] \(\frac{a}{3}\)
    \item[C)] \(\frac{a}{4}\)
    \item[D)] \(\frac{\sqrt{3}a}{2}\)
\end{enumerate}

\vspace{0.cm} % Пустое место для решения

\item Развертка правильной шестиугольной призмы показана на рисунке. Если \(AB = 8\) м, \(CD = 6\) м, найдите объем правильной шестиугольной призмы. \\
\begin{center}
    \centering
    \includegraphics[width=0.5\linewidth]{12.png}
\end{center}
\textbf{}
\begin{enumerate}
    \item[A)] \(120\sqrt{3}\)
    \item[B)] \(108\sqrt{3}\)
    \item[C)] \(144\sqrt{3}\)
    \item[D)] \(96\sqrt{3}\)
\end{enumerate}

\vspace{0.cm} % Пустое место для решения

\item Сумма площадей двух оснований цилиндра равна площади его боковой поверхности. Если его объем равен \(64\pi\), найдите площадь осевого сечения цилиндра. \\
\textbf{}
\begin{enumerate}
    \item[A)] 32
    \item[B)] 16
    \item[C)] 8
    \item[D)] 20
\end{enumerate}

\vspace{0.cm} % Пустое место для решения

\item В шар вписана правильная четырехугольная пирамида. Если высота пирамиды делится экватором шара на отрезки длиной 5 и 4, найдите объем пирамиды. \\
\textbf{}
\begin{enumerate}
    \item[A)] 32
    \item[B)] 54
    \item[C)] 64
    \item[D)] 81
\end{enumerate}

\vspace{0.cm} % Пустое место для решения

\item В правильной треугольной пирамиде \(SABC\) длина бокового ребра в 2 раза больше стороны основания. В треугольнике \(SAB\) проведена высота \(AH\), а в треугольнике \(ABC\) — медиана \(BM\). Найдите отношение длины отрезка \(AH\) к длине отрезка \(BM\). \\
\textbf{}
\begin{enumerate}
    \item[A)] 2:1
    \item[B)] \(\frac{\sqrt{5}}{2}\)
    \item[C)] \(\frac{\sqrt{3}}{4}\)
    \item[D)] \(\frac{13}{5}\)
\end{enumerate}

\vspace{0.cm} % Пустое место для решения

\item Объем правильной треугольной призмы равен 16. При какой длине стороны основания полная поверхность призмы будет наименьшей? \\
\textbf{}
\begin{enumerate}
    \item[A)] 3
    \item[B)] 4
    \item[C)] 2
    \item[D)] \(3\sqrt{2}\)
\end{enumerate}

\vspace{0.cm} % Пустое место для решения
\item В правильной усечённой треугольной пирамиде стороны оснований равны 3 и
7. В неё вписан шар. Найдите площадь боковой поверхности пирамиды.
\begin{enumerate}
    \item[A)] \(25\sqrt{3}\)
    \item[B)] \(20\sqrt{3}\)
    \item[C)] \(15\sqrt{5}\)
    \item[D)] 27
\end{enumerate}

\end{enumerate}

\end{document}